
d'une
structure de \(\lambda\)-anneau, les relations (2.4.2) passant au quotient.

Toute intersection de \(\lambda\)-idéaux est un \(\lambda\)-idéal. En particulier, il existe un plus petit \(\lambda\)-idéal contenant une famille donnée d'éléments \((x_i)_{i\in I}\)~; d'après les relations (2.4.2), c'est l'idéal (ordinaire) engendré par les \(\lambda^n(x_i)\), \((n\ge1)\).

Soient toujours \(K\) un \(\lambda\)-anneau, \((X_i)_{i\in I}\) une famille d'indéterminées, \(A=K[(X_i),(\lambda^n X_i)]\), et soit \((P_{i'})_{i'\in I'}\) une famille d'éléments de \(A\) : soit \(J\) le \(\lambda\)-idéal engendré par les \(P_{i'}\).

\textbf{Définition 4.6.} \uline{Le} \(\lambda\)-\uline{anneau} \(A/J\) \uline{sera appelé} \(\lambda\)-anneau engendré sur \(K\) par la famille de générateurs \((X_i)_{i\in I}\) soumis aux relations \((P_{i'})_{i'\in I'}\).

On peut également se donner des relations entre les \(X_i\), \(\gamma^n X_i\) et faire la même construction, en remarquant qu'un idéal est un \(\lambda\)-idéal si et seulement si il est stable pour les opérations \(\gamma^n\), \(n\ge1\), comme il résulte des relations 3.2.

Soit \(B\) une \(K\)-\(\lambda\)-algèbre, \((b_i)_{i\in I}\) une famille d'éléments de \(B\). Les \(P_{i'}\) étant des polynômes par rapport aux indéterminées \(X_i\), \(\lambda^n X_i\), on peut substituer à ces indéterminées les \(b_i\), \(\lambda^n(b_i)\)~; si l'élément obtenu est nul, on dira que les \(b_i\) satisfont les relations \(P_{i'}\). La \(K\)-\(\lambda\)-algèbre \(A/J\) est alors caractérisée par la propriété universelle suivante :

\textbf{Proposition 4.7.} \uline{Pour toute} \(K\)-\(\lambda\)-\uline{algèbre} \(B\) \uline{et toute famille} \((b_i)_{i\in I}\) \uline{d'éléments de} \(B\), \uline{satisfaisant les relations} \((P_{i'})_{i'\in I'}\), \uline{il existe un} \(K\)-\(\lambda\)-\uline{homomorphisme et un seul} \(f':A/J\to B\) \uline{tel que pour tout} \(i\), on ait \(f'(X_i)=b_i\).

En effet, le noyau de l'homomorphisme \(f\) défini en 4.2 contient les \(P_{i'}\), donc \(J\) puisque c'est un \(\lambda\)-idéal~; \(f\) passe donc au quotient et définit 
